\section{IMUの計測値に対する残差に関するヤコビアンの導出}

式(\ref{eq:imu_residual})に示すIMUの計測値に対する残差を計算するにあたって使用される角速度$\boldsymbol \omega(t)$と加速度${\bf a}(t)$は，制御点を用いてそれぞれ以下のように計算される．
%
\begin{align}
  \boldsymbol \omega(t)
  %
  = \left( R(t)^{\top} \dot{ R }(t) \right)^{\vee}
  %
  = \left( R(t)^{\top} R_{0} \left( \dot{ A_{1} } A_{2} A_{3} + A_{1} \dot{ A_{2} } A_{3} + A_{1} A_{2} \dot{ A_{3} } \right) \right)^{\vee}
  \label{eq:continuous_time_angular_velocity}
\end{align}
%
\begin{align}
  {\bf a}(t)
  =
  \ddot{ {\bf p} }(t)
  =
  \sum_{j=1}^{3} \ddot{ \lambda }_{j}(t) {\bf d}_{j}
  \label{eq:continuous_time_acceleration}
\end{align}
%
なお式(\ref{eq:imu_residual})，(\ref{eq:continuous_time_angular_velocity})，(\ref{eq:continuous_time_acceleration})から明らかなように，角速度に関する残差${}^{I}{\bf r}_{\omega}(t)$は姿勢$R_{j}$，加速度に関する残差${}^{I}{\bf r}_{a}(t)$は姿勢$R_{j}$と並進${\bf p}_{j}$に影響を受ける．
そのため，求めるヤコビアン$J$は以下となる．
%
\begin{align}
  \begin{split}
  J
  & = 
  \left( \begin{matrix}
    \frac{ \partial {}^{I}{\bf r}_{\omega}(t) }{ \partial R_{0} } &
    \frac{ \partial {}^{I}{\bf r}_{\omega}(t) }{ \partial {\bf p}_{0} } &
    \frac{ \partial {}^{I}{\bf r}_{\omega}(t) }{ \partial R_{1} } &
    \frac{ \partial {}^{I}{\bf r}_{\omega}(t) }{ \partial {\bf p}_{1} } &
    \frac{ \partial {}^{I}{\bf r}_{\omega}(t) }{ \partial R_{2} } &
    \frac{ \partial {}^{I}{\bf r}_{\omega}(t) }{ \partial {\bf p}_{2} } &
    \frac{ \partial {}^{I}{\bf r}_{\omega}(t) }{ \partial R_{3} } &
    \frac{ \partial {}^{I}{\bf r}_{\omega}(t) }{ \partial {\bf p}_{3} } &
    \frac{ \partial {}^{I}{\bf r}_{\omega}(t) }{ \partial {\bf b}_{g}^{k+1} } &
    \frac{ \partial {}^{I}{\bf r}_{\omega}(t) }{ \partial {\bf b}_{a}^{k+1} } \\
    %
    \frac{ \partial {}^{I}{\bf r}_{a}(t) }{ \partial R_{0} } &
    \frac{ \partial {}^{I}{\bf r}_{a}(t) }{ \partial {\bf p}_{0} } &
    \frac{ \partial {}^{I}{\bf r}_{a}(t) }{ \partial R_{1} } &
    \frac{ \partial {}^{I}{\bf r}_{a}(t) }{ \partial {\bf p}_{1} } &
    \frac{ \partial {}^{I}{\bf r}_{a}(t) }{ \partial R_{2} } &
    \frac{ \partial {}^{I}{\bf r}_{a}(t) }{ \partial {\bf p}_{2} } &
    \frac{ \partial {}^{I}{\bf r}_{a}(t) }{ \partial R_{3} } &
    \frac{ \partial {}^{I}{\bf r}_{a}(t) }{ \partial {\bf p}_{3} } &
    \frac{ \partial {}^{I}{\bf r}_{a}(t) }{ \partial {\bf b}_{g}^{k+1} } &
    \frac{ \partial {}^{I}{\bf r}_{a}(t) }{ \partial {\bf b}_{a}^{k+1} } \\
  \end{matrix} \right) \\
  %
  & = 
  \left( \begin{matrix}
    \frac{ \partial {}^{I}{\bf r}_{\omega}(t) }{ \partial R_{0} } &
    0 &
    \frac{ \partial {}^{I}{\bf r}_{\omega}(t) }{ \partial R_{1} } &
    0 &
    \frac{ \partial {}^{I}{\bf r}_{\omega}(t) }{ \partial R_{2} } &
    0 &
    \frac{ \partial {}^{I}{\bf r}_{\omega}(t) }{ \partial R_{3} } &
    0 &
    I &
    0 \\
    %
    \frac{ \partial {}^{I}{\bf r}_{a}(t) }{ \partial R_{0} } &
    \frac{ \partial {}^{I}{\bf r}_{a}(t) }{ \partial {\bf p}_{0} } &
    \frac{ \partial {}^{I}{\bf r}_{a}(t) }{ \partial R_{1} } &
    \frac{ \partial {}^{I}{\bf r}_{a}(t) }{ \partial {\bf p}_{1} } &
    \frac{ \partial {}^{I}{\bf r}_{a}(t) }{ \partial R_{2} } &
    \frac{ \partial {}^{I}{\bf r}_{a}(t) }{ \partial {\bf p}_{2} } &
    \frac{ \partial {}^{I}{\bf r}_{a}(t) }{ \partial R_{3} } &
    \frac{ \partial {}^{I}{\bf r}_{a}(t) }{ \partial {\bf p}_{3} } &
    0 &
    I \\
  \end{matrix} \right) \in \mathbb{R}^{6 \times 30}
  \end{split}
\end{align}
%
なおIMUの計測値バイアスに対するヤコビアンも記述しているが，これらの計算結果は式(\ref{eq:imu_residual})より明らかであるため省略する．
以下ではまず，加速度に関するヤコビアンの導出から考える．





\subsection{加速度残差に関するヤコビアンの導出}

\subsubsection{並進に関するヤコビアンの導出}

まず加速度に関する残差を並進成分に依存する変数のみで以下の様に記述する．
%
\begin{align}
  {}^{I}{\bf r}_{a}
  \left(
    {\bf d}_{1} \left( {\bf p}_{0}, {\bf p}_{1} \right),
    {\bf d}_{2} \left( {\bf p}_{1}, {\bf p}_{2} \right),
    {\bf d}_{3} \left( {\bf p}_{2}, {\bf p}_{3} \right)
  \right)
\end{align}
%
これより，$j=1,2$の場合は連鎖則を用いて以下のようにヤコビアンが計算できる．
%
\begin{align}
  \frac{ \partial {}^{I}{\bf r}_{a}(t) }{ \partial {\bf p}_{j} }
  =
  \frac{ \partial {}^{I}{\bf r}_{a}(t) }{ \partial {\bf d}_{j} }
  \frac{ \partial {\bf d}_{j} }{ \partial {\bf p}_{j} }
  +
  \frac{ \partial {}^{I}{\bf r}_{a}(t) }{ \partial {\bf d}_{j+1} }
  \frac{ \partial {\bf d}_{j+1} }{ \partial {\bf p}_{j} }
\end{align}
%
$j=0$と$j=3$の場合は以下となる．
%
\begin{align}
  \frac{ \partial {}^{I}{\bf r}_{a}(t) }{ \partial {\bf p}_{0} }
  =
  \frac{ \partial {}^{I}{\bf r}_{a}(t) }{ \partial {\bf d}_{1} }
  \frac{ \partial {\bf d}_{1} }{ \partial {\bf p}_{0} }
\end{align}
%
\begin{align}
  \frac{ \partial {}^{I}{\bf r}_{a}(t) }{ \partial {\bf p}_{3} }
  =
  \frac{ \partial {}^{I}{\bf r}_{a}(t) }{ \partial {\bf d}_{3} }
  \frac{ \partial {\bf d}_{3} }{ \partial {\bf p}_{3} }
\end{align}

ここでまず，式(\ref{eq:continuous_time_acceleration})を用いて加速度に関する残差を以下の様に書き直す．
%
\begin{align}
  {}^{I}{\bf r}_{a}(t) = R(t)^{\top} \left( \sum_{j=1}^{3} \ddot{ \lambda }_{j}(t) {\bf d}_{j} + {\bf g} \right) - {}^{I}{\bf a} + {\bf b}_{a}^{k+1}
\end{align}
%
これから，$\partial {}^{I}{\bf r}_{a}(t) / \partial {\bf d}_{j}$は明らかに以下となる．
%
\begin{align}
  \frac{ \partial {}^{I}{\bf r}_{a}(t) }{ \partial {\bf d}_{j} } = R(t)^{\top} \ddot{ \lambda }_{j}(t)
\end{align}
%
なお${\bf d}_{j} = {\bf p}_{j} - {\bf p}_{j-1}$であるため，$\partial {\bf d}_{j} / \partial {\bf p}_{j} = I$，$\partial {\bf d}_{j+1} / \partial {\bf p}_{j} = -I$は明らかである．
よって，加速度の残差に関する並進のヤコビアンとしてそれぞれ以下を得る．
%
\begin{align}
  \begin{gathered}
    \frac{ {}^{I}{\bf r}_{a}(t) }{ \partial {\bf p}_{0} }
    =
    -\ddot{ \lambda }_{1} R(t)^{\top} \\
    %
    \frac{ {}^{I}{\bf r}_{a}(t) }{ \partial {\bf p}_{1} }
    =
    \left( \ddot{ \lambda }_{1} - \ddot{ \lambda }_{2} \right) R(t)^{\top} \\
    %
    \frac{ {}^{I}{\bf r}_{a}(t) }{ \partial {\bf p}_{2} }
    =
    \left( \ddot{ \lambda }_{2} - \ddot{ \lambda }_{3} \right) R(t)^{\top} \\
    %
    \frac{ {}^{I}{\bf r}_{a}(t) }{ \partial {\bf p}_{3} }
    =
    \ddot{ \lambda }_{3} R(t)^{\top} \\
    %
  \end{gathered}
\end{align}








\subsubsection{回転に関するヤコビアンの導出}

まず加速度に関する残差をまず姿勢$R_{j}$にのみ依存する関数として記述する．
%
\begin{align}
  {}^{I}{\bf r}_{a}
  \left(
    R
    \left(
      R_{0},
      \boldsymbol \phi_{1} \left( R_{0}, R_{1} \right),
      \boldsymbol \phi_{2} \left( R_{1}, R_{2} \right),
      \boldsymbol \phi_{3} \left( R_{2}, R_{3} \right)
    \right)
  \right)
\end{align}
%
ここから，$j=1,2$の場合のヤコビアンは連鎖則により以下のように計算できる．
%
\begin{align}
  \frac{ \partial {}^{I}{\bf r}_{a}(t) }{ \partial R_{j} }
  =
  \frac{ \partial {}^{I}{\bf r}_{a}(t) }{ \partial R(t) }  
  \frac{ \partial R(t) }{ \partial \boldsymbol \phi_{j} }  
  \frac{ \partial \boldsymbol \phi_{j} }{ \partial R_{j} }  
  +
  \frac{ \partial {}^{I}{\bf r}_{a}(t) }{ \partial R(t) }  
  \frac{ \partial R(t) }{ \partial \boldsymbol \phi_{j+1} }  
  \frac{ \partial \boldsymbol \phi_{j+1} }{ \partial R_{j} }  
\end{align}
%
$j=0$と$j=3$の場合は以下となる．
%
\begin{align}
  \frac{ \partial {}^{I}{\bf r}_{a}(t) }{ \partial R_{0} }
  =
  \frac{ \partial {}^{I}{\bf r}_{a}(t) }{ \partial R(t) }  
  \frac{ \partial R(t) }{ \partial R_{0} }  
  +
  \frac{ \partial {}^{I}{\bf r}_{a}(t) }{ \partial R(t) }  
  \frac{ \partial R(t) }{ \partial \boldsymbol \phi_{j} }  
  \frac{ \partial \boldsymbol \phi_{j} }{ \partial R_{j} }  
\end{align}
%
\begin{align}
  \frac{ \partial {}^{I}{\bf r}_{a}(t) }{ \partial R_{3} }
  =
  \frac{ \partial {}^{I}{\bf r}_{a}(t) }{ \partial R(t) }  
  \frac{ \partial R(t) }{ \partial \boldsymbol \phi_{3} }  
  \frac{ \partial \boldsymbol \phi_{3} }{ \partial R_{3} }  
\end{align}
%
ここで$\partial {}^{I}{\bf r}_{a}(t) / \partial R(t)$以外の項は，スキャンマッチングに関する残差のヤコビアンを計算する際にも利用されているため，詳細は\ref{sec:scan_matching}章に譲る．

$\partial {}^{I}{\bf r}_{a}(t) / \partial R(t)$を計算するにあたり，以下の摂動を考える．
%
\begin{align}
  \begin{split}
    & \left( {\rm Exp} \left( \delta \boldsymbol \theta \right) R(t) \right)^{\top} \left( {\bf a}(t) + {\bf g} \right) - {}^{I}{\bf a} + {\bf b}_{a}^{k+1} - \left( R(t)^{\top}  \left( {\bf a}(t) + {\bf g} \right) - {}^{I}{\bf a} + {\bf b}_{a}^{k+1} \right) \\
    %
    = & R(t)^{\top}  \left( {\rm Exp} \left( -\delta \boldsymbol \theta \right) - I \right) \left( {\bf a}(t) + {\bf g} \right) \\
    %
    = & R(t)^{\top}  \left( -\delta \boldsymbol \theta\right)^{\wedge} \left( {\bf a}(t) + {\bf g} \right) \\
    %
    = & R(t)^{\top} \left( {\bf a}(t) + {\bf g} \right)^{\wedge} \delta \boldsymbol \theta
  \end{split}
\end{align}
%
よって以下を得る．
%
\begin{align}
  \frac{ \partial {}^{I}{\bf r}_{a}(t) }{ \partial R(t) } = R(t)^{\top}  \left( {\bf a}(t) + {\bf g} \right)^{\wedge}
\end{align}












\subsection{角速度残差に関するヤコビアンの導出}

角速度の残差に関するヤコビアンの導出をするにあたり，まず角速度$\boldsymbol \omega(t)$を算出する際の歪対象行列$\Omega(t)$を以下の様に式変形する．
%
\begin{align}
  \begin{split}
    \Omega(t) = & R(t)^{\top}  R_{0} \left( \dot{ A_{1} } A_{2} A_{3} + A_{1} \dot{ A_{2} } A_{3} + A_{1} A_{2} \dot{ A_{3} } \right) \\
    %
    = & A_{3}^{\top} A_{2}^{\top} A_{1}^{\top} R_{0}^{\top} R_{0} \left( \left( \dot{ \lambda }_{1} \boldsymbol \phi_{1} \right)^{\wedge} A_{1} A_{2} A_{3} + A_{1} \left( \dot{ \lambda }_{2} \boldsymbol \phi_{2} \right)^{\wedge} A_{2} A_{3} + A_{1} A_{2} \left( \dot{ \lambda }_{3} \boldsymbol \phi_{3} \right)^{\wedge} A_{3} \right) \\
    %
    = & \left( A_{3}^{\top} A_{2}^{\top} A_{1}^{\top} \dot{ \lambda }_{1} \boldsymbol \phi_{1} \right)^{\wedge} + \left( A_{3}^{\top} A_{2}^{\top} \dot{ \lambda }_{2} \boldsymbol \phi_{2} \right)^{\wedge} + \left( A_{3}^{\top} \dot{ \lambda }_{3} \boldsymbol \phi_{3} \right)^{\wedge} \\
    %
    = & \left( A_{3}^{\top} A_{2}^{\top} A_{1}^{\top} \dot{ \lambda }_{1} \boldsymbol \phi_{1} + A_{3}^{\top} A_{2}^{\top} \dot{ \lambda }_{2} \boldsymbol \phi_{2} + A_{3}^{\top} \dot{ \lambda }_{3} \boldsymbol \phi_{3} \right)^{\wedge}
  \end{split}
\end{align}
%
なお，以下の関係を用いている．
%
\begin{align}
  \dot{ A }_{j} = \left( \dot{ \lambda }_{j} \boldsymbol \phi_{j} \right)^{\wedge} A_{j}
\end{align}
%
これより，角速度$\boldsymbol \omega(t)$は以下のように変形できる．
%
\begin{align}
  \boldsymbol \omega(t) = A_{3}^{\top} A_{2}^{\top} A_{1}^{\top} \dot{ \lambda }_{1} \boldsymbol \phi_{1} + A_{3}^{\top} A_{2}^{\top} \dot{ \lambda }_{2} \boldsymbol \phi_{2} + A_{3}^{\top} \dot{ \lambda }_{3} \boldsymbol \phi_{3}
  \label{eq:angular_velocity_mod}
\end{align}
%

式(\ref{eq:angular_velocity_mod})に示した角速度$\boldsymbol \omega(t)$を参照し，角速度に関する残差を以下の様に記述する．
%
\begin{align}
  {}^{I}{\bf r}_{\omega}
  \left(
    \boldsymbol \omega
    \left(
      \boldsymbol \phi_{1} \left( R_{0}, R_{1} \right),
      \boldsymbol \phi_{2} \left( R_{1}, R_{2} \right),
      \boldsymbol \phi_{3} \left( R_{2}, R_{3} \right)
    \right)
  \right)
\end{align}
%
これより，$j=1,2$の場合の$R_{j}$に関するヤコビアンは，連鎖則を用いて以下の様に計算できる．
%
\begin{align}
  \frac{ \partial {}^{I}{\bf r}_{\omega} }{ \partial R_{j} }
  =
  \frac{ \partial {}^{I}{\bf r}_{\omega} }{ \partial \boldsymbol \omega(t) }
  \frac{ \partial \omega(t) }{ \partial \boldsymbol \phi_{j} }
  \frac{ \partial \boldsymbol \phi_{j} }{ \partial R_{j} }
  +
  \frac{ \partial {}^{I}{\bf r}_{\omega} }{ \partial \boldsymbol \omega(t) }
  \frac{ \partial \omega(t) }{ \partial \boldsymbol \phi_{j+1} }
  \frac{ \partial \boldsymbol \phi_{j+1} }{ \partial R_{j} }
\end{align}
%
また$j=0$，および$j=3$の場合は以下となる．
%
\begin{align}
  \frac{ \partial {}^{I}{\bf r}_{\omega} }{ \partial R_{0} }
  =
  \frac{ \partial {\bf r}_{\omega} }{ \partial \boldsymbol \omega(t) }
  \frac{ \partial \boldsymbol \omega(t) }{ \partial \boldsymbol \phi_{1} }
  \frac{ \partial \boldsymbol \phi_{1} }{ \partial R_{0} }
\end{align}
%
\begin{align}
  \frac{ \partial {}^{I}{\bf r}_{\omega} }{ \partial R_{3} }
  =
  \frac{ \partial {\bf r}_{\omega} }{ \partial \boldsymbol \omega(t) }
  \frac{ \partial \boldsymbol \omega(t) }{ \partial \boldsymbol \phi_{3} }
  \frac{ \partial \boldsymbol \phi_{3} }{ \partial R_{3} }
\end{align}
%
なお$\boldsymbol \omega(t)$は$R_{0}$に依存しないため，$\partial \boldsymbol \omega(t) / \partial R_{0}$はゼロとなる．
また式(\ref{eq:imu_residual})より，明らかに$\partial {}^{I}{\bf r}_{\omega}(t) / \partial \boldsymbol \omega(t) = I$であり，$\partial \boldsymbol \phi_{j} / \partial R_{j}$と$\partial \boldsymbol \phi_{j+1} / \partial R_{j}$はスキャンマッチングに関するヤコビアンを計算する際にも導出しているため詳細は省略する．

最後に$\partial \boldsymbol \omega(t) / \partial \boldsymbol \phi_{j}$を考える．
これは，$\boldsymbol \omega(t)$を$\boldsymbol \phi_{j}$に関する関数とみなし，$\boldsymbol \omega \left( \boldsymbol \phi_{j} + \delta \boldsymbol \phi_{j} \right) - \boldsymbol \omega \left( \boldsymbol \phi_{j} \right) = \boldsymbol \omega \left( \boldsymbol \phi_{j} \right) + J \delta \boldsymbol \phi_{j} + O \left( \left\| \delta \boldsymbol \phi_{j} \right\|^{2} \right)$となるヤコビアン$J$を求めることで計算できる．
例えば$\partial \boldsymbol \omega(t) / \partial \boldsymbol \phi_{1}$に関しては，以下の様に計算できる．
%
\begin{align}
  \begin{split}
    & A_{3}^{\top} A_{2}^{\top} \left( {\rm Exp} \left( \lambda_{1} \left( \boldsymbol \phi_{1} + \delta \boldsymbol \phi_{1} \right) \right) \right)^{\top} \dot{ \lambda }_{1} \left( \boldsymbol \phi_{1} + \delta \boldsymbol \phi_{1} \right) + A_{3}^{\top} A_{2}^{\top} \dot{ \lambda }_{2} \boldsymbol \phi_{2} + A_{3}^{\top} \dot{ \lambda }_{3} \boldsymbol \phi_{3} \\
    & - A_{3}^{\top} A_{2}^{\top} A_{1}^{\top} \dot{ \lambda }_{1} \boldsymbol \phi_{1} + A_{3}^{\top} A_{2}^{\top} \dot{ \lambda }_{2} \boldsymbol \phi_{2} + A_{3}^{\top} \dot{ \lambda }_{3} \boldsymbol \phi_{3} \\
    %
    \simeq & A_{3}^{\top} A_{2}^{\top} \left( {\rm Exp} \left( J_{l} \left( \lambda_{1} \boldsymbol \phi_{1} \right) \lambda_{1} \delta \boldsymbol \phi_{1} \right) {\rm Exp} \left( \lambda_{1} \boldsymbol \phi_{1} \right) \right)^{\top} \dot{ \lambda }_{1} \left( \boldsymbol \phi_{1} + \delta \boldsymbol \phi_{1} \right)
    - A_{3}^{\top} A_{2}^{\top} A_{1}^{\top} \dot{ \lambda }_{1} \boldsymbol \phi_{1} \\
    %
    = & A_{3}^{\top} A_{2}^{\top} {\rm Exp} \left( -\lambda_{1} \boldsymbol \phi_{1} \right) {\rm Exp} \left( -J_{l} \left( \lambda_{1} \boldsymbol \phi_{1} \right) \lambda_{1} \delta \boldsymbol \phi_{1} \right) \dot{ \lambda }_{1} \left( \boldsymbol \phi_{1} + \delta \boldsymbol \phi_{1} \right)
    - A_{3}^{\top} A_{2}^{\top} A_{1}^{\top} \dot{ \lambda }_{1} \boldsymbol \phi_{1} \\
    %
    = & A_{3}^{\top} A_{2}^{\top} A_{1}^{\top} \left( {\rm Exp} \left( -J_{l} \left( \lambda_{1} \boldsymbol \phi_{1} \right) \lambda_{1} \delta \boldsymbol \phi_{1} \right) \dot{ \lambda }_{1} \left( \boldsymbol \phi_{1} + \delta \boldsymbol \phi_{1} \right) - \dot{ \lambda }_{1} \boldsymbol \phi_{1} \right) \\
    %
    = & A_{3}^{\top} A_{2}^{\top} A_{1}^{\top} \left( \left( {\rm Exp} \left( -J_{l} \left( \lambda_{1} \boldsymbol \phi_{1} \right) \lambda_{1} \delta \boldsymbol \phi_{1} \right) - I \right) \dot{ \lambda}_{1} \boldsymbol \phi_{1} + {\rm Exp} \left( -J_{l} \left( \lambda_{1} \boldsymbol \phi_{1} \right) \lambda_{1} \delta \boldsymbol \phi_{1} \right) \dot{ \lambda }_{1} \delta \boldsymbol \phi_{1} \right) \\
    %
    \simeq & A_{3}^{\top} A_{2}^{\top} A_{1}^{\top} \left( \left( I - \left( J_{l} \left( \lambda_{1} \boldsymbol \phi_{1} \right) \lambda_{1} \delta \boldsymbol \phi_{1} \right)^{\wedge} + I \right) \dot{ \lambda }_{1} \boldsymbol \phi_{1} + \left( I - \left( J_{l} \left( \lambda_{1} \boldsymbol \phi_{1} \right) \lambda_{1} \delta \boldsymbol \phi_{1}  \right)^{\wedge} \right) \dot{ \lambda }_{1} \delta \boldsymbol \phi_{1} \right) \\
    %
    \simeq & A_{3}^{\top} A_{2}^{\top} A_{1}^{\top} \left( -\left( J_{l} \left( \lambda_{1} \boldsymbol \phi_{1} \right) \lambda_{1} \delta \boldsymbol \phi_{1} \right)^{\wedge} \dot{ \lambda }_{1} \boldsymbol \phi_{1} + \dot{ \lambda }_{1} \delta \boldsymbol \phi_{1} \right) \\
    %
    = & A_{3}^{\top} A_{2}^{\top} A_{1}^{\top} \left( \left( \dot{ \lambda }_{1} \boldsymbol \phi_{1} \right)^{\wedge} J_{l} \left( \lambda_{1} \boldsymbol \phi_{1} \right) \lambda_{1} \delta \boldsymbol \phi_{1} + \dot{ \lambda }_{1} \delta \boldsymbol \phi_{1} \right) \\
    %
    = & A_{3}^{\top} A_{2}^{\top} A_{1}^{\top} \left( \left( \dot{ \lambda }_{1} \boldsymbol \phi_{1} \right)^{\wedge} J_{l} \left( \lambda_{1} \boldsymbol \phi_{1} \right) \lambda_{1} + \dot{ \lambda }_{1} I \right) \delta \boldsymbol \phi_{1} \\
  \end{split}
\end{align}
%
すなわち，$\partial \boldsymbol \omega(t) / \partial \boldsymbol \phi_{1}$は以下となる．
%
\begin{align}
  \frac{ \partial \boldsymbol \omega(t) }{ \partial \boldsymbol \phi_{1} }
  =
  A_{3}^{\top} A_{2}^{\top} A_{1}^{\top} \left( \left( \dot{ \lambda }_{1} \boldsymbol \phi_{1} \right)^{\wedge} J_{l} \left( \lambda_{1} \boldsymbol \phi_{1} \right) \lambda_{1} + \dot{ \lambda }_{1} I \right)
\end{align}
%
なお計算の途中で${\rm Exp} \left( -\lambda_{1} \boldsymbol \phi_{1} \right) = A_{1}^{\top}$，$\left( J_{l} \left( \lambda_{1} \boldsymbol \phi_{1} \right) \lambda_{1} \delta \boldsymbol \phi_{1}  \right)^{\wedge} \dot{ \lambda }_{1} \delta \boldsymbol \phi_{1} \simeq {\bf 0}$を利用している．
$\partial \boldsymbol \omega(t) / \partial \boldsymbol \phi_{2}$と$\partial \boldsymbol \omega(t) / \partial \boldsymbol \phi_{3}$も同様に計算できる．
しかし計算が煩雑となるため詳細は省くが，計算結果は以下となる．
%
\begin{align}
  \frac{ \partial \boldsymbol \omega(t) }{ \partial \boldsymbol \phi_{2} }
  =
  A_{3}^{\top} A_{2}^{\top}
  \left(
    \left( A_{1}^{\top} \dot{ \lambda }_{1} \boldsymbol \phi_{1} \right)^{\wedge} J_{l} \left( \lambda_{2} \boldsymbol \phi_{2} \right) \lambda_{2}
    + \left( \dot{ \lambda }_{2} \boldsymbol \phi_{2} \right)^{\wedge} J_{l} \left( \lambda_{2} \boldsymbol \phi_{2} \right) \lambda_{2}
    + \dot{ \lambda }_{2} I
  \right)
\end{align}
%
\begin{align}
  \frac{ \partial \boldsymbol \omega(t) }{ \partial \boldsymbol \phi_{3} }
  =
  A_{3}^{\top}
  \left(
    \left( A_{2}^{\top} A_{1}^{\top} \dot{ \lambda }_{1} \boldsymbol \phi_{1} \right)^{\wedge} J_{l} \left( \lambda_{3} \boldsymbol \phi_{3} \right) \lambda_{3}
    + \left( A_{2}^{\top} \dot{ \lambda }_{2} \boldsymbol \phi_{2} \right)^{\wedge} J_{l} \left( \lambda_{3} \boldsymbol \phi_{3} \right) \lambda_{3}
    + \left( \dot{ \lambda }_{3} \boldsymbol \phi_{3} \right)^{\wedge} J_{l} \left( \lambda_{3} \boldsymbol \phi_{3} \right) \lambda_{3}
    + \dot{ \lambda }_{3} I
  \right)
\end{align}
%
以上の計算結果から，IMUの角速度の計測値に対する残差に関するヤコビアンを導出することができる．

